\section*{Capítulo \ref{ch:g12:comb} --- Combinatoria y conteo}

\begin{solution}{exo:g12:comb:1}
Principio del producto:
$26^2 \times 10^3 \times 26^2 = 26^4 \times 10^3 = 456\,976\,000$.

Sin repetir caracteres, las cuatro letras han de ser distintas
($26 \times 25 \times 24 \times 23$ maneras, rellenando las posiciones de
las letras en orden) y las tres cifras también
($10 \times 9 \times 8$):
\[
26 \times 25 \times 24 \times 23 \times 10 \times 9 \times 8
= 358\,800 \times 720 = 258\,336\,000 .
\]
\end{solution}

\begin{solution}{exo:g12:comb:2}
$\dbinom83 = \dfrac{8 \times 7 \times 6}{3!} = 56$;
$\dbinom{10}{8} = \dbinom{10}{2} = \dfrac{10 \times 9}{2} = 45$;
\[
\frac{\binom n2}{\binom{n+1}2}
= \frac{n(n-1)/2}{(n+1)n/2} = \frac{n-1}{n+1}.
\]
\end{solution}

\begin{solution}{exo:g12:comb:3}
Elegimos la comisión: $\binom{30}{4}$ maneras. Elegimos después el
presidente y el tesorero entre los 4, en orden: $4 \times 3 = 12$ maneras.
Total:
\[
\binom{30}{4} \times 12 = 27\,405 \times 12 = 328\,860 .
\]
\end{solution}

\begin{solution}{exo:g12:comb:4}
\[
(x+2)^5 = x^5 + 10x^4 + 40x^3 + 80x^2 + 80x + 32 ,
\]
\[
(1-x)^6 = 1 - 6x + 15x^2 - 20x^3 + 15x^4 - 6x^5 + x^6 .
\]
En $(2x+3)^7$, el término en $x^3$ es
$\binom{7}{3}(2x)^3\,3^4 = 35 \times 8 \times 81\, x^3$: el coeficiente es
$22\,680$.
\end{solution}

\begin{solution}{exo:g12:comb:5}
\emph{1.} $\dbinom{52}{5} = 2\,598\,960$.

\emph{2.} Exactamente un as: lo elegimos ($4$ maneras) y completamos con
$4$ cartas que no sean ases:
$4 \times \binom{48}{4} = 4 \times 194\,580 = 778\,320$. Al menos un as:
por el \omterm{def:g10:proba:operations}{suceso contrario},
$\binom{52}{5} - \binom{48}{5} = 2\,598\,960 - 1\,712\,304 = 886\,656$.

\emph{3.} Elegimos el valor del trío ($13$), sus palos
($\binom43 = 4$), el valor de la pareja ($12$ restantes) y sus palos
($\binom42 = 6$): $13 \times 4 \times 12 \times 6 = 3744$.
\end{solution}

\begin{solution}{exo:g12:comb:6}
ROMA tiene 4 letras distintas: $4! = 24$ anagramas.

BANANA tiene 6 letras: tres aes, dos enes y una be. Elegimos las posiciones
de las aes ($\binom63$) y después las de las enes entre las restantes
($\binom32$); la be ocupa el último hueco:
\[
\binom{6}{3}\binom{3}{2} = 20 \times 3 = 60 .
\]
(Equivalentemente, $\frac{6!}{3!\,2!\,1!} = 60$.)
\end{solution}

\begin{solution}{exo:g12:comb:7}
\emph{Algebraicamente:}
\[
k\binom nk = \frac{k\,n!}{k!(n-k)!} = \frac{n!}{(k-1)!\,(n-k)!}
= n\,\frac{(n-1)!}{(k-1)!\bigl((n-1)-(k-1)\bigr)!} = n\binom{n-1}{k-1}.
\]
\emph{Por doble conteo:} contamos las parejas (comisión de $k$ personas,
presidente de la comisión). O bien elegimos la comisión ($\binom nk$) y
después su presidente ($k$): $k\binom nk$ parejas. O bien elegimos primero
al presidente ($n$ opciones) y después a los otros $k-1$ miembros entre los
$n-1$ restantes: $n\binom{n-1}{k-1}$ parejas.
\end{solution}

\begin{solution}{exo:g12:comb:8}
Un camino consta de exactamente $m + n$ pasos, de los cuales $m$ van hacia
el este y $n$ hacia el norte; queda completamente determinado por el
conjunto de instantes (entre los $m+n$) en los que se da un paso hacia el
este. Hay $\binom{m+n}{m}$ elecciones así.
\end{solution}

\begin{solution}{exo:g12:comb:9}
Repartimos un conjunto de $m + n$ personas en un grupo $A$ de $m$ y un
grupo $B$ de $n$. Un subconjunto de $k$ elementos contiene un cierto número
$j$ de miembros de $A$ ($0 \leq j \leq k$) y $k - j$ de $B$; para $j$ fijo
hay $\binom mj \binom{n}{k-j}$ subconjuntos así, y el principio de la suma
sobre $j$ da la identidad de Vandermonde.

Con $m = n = k$:
\[
\binom{2n}{n} = \sum_{j=0}^n \binom nj \binom{n}{n-j}
= \sum_{j=0}^n \binom nj^{2},
\]
usando la simetría $\binom{n}{n-j} = \binom nj$.
\end{solution}

\begin{solution}{exo:g12:comb:10}
\emph{Con el \cref{exo:g12:comb:7}:}
\[
\sum_{k=1}^n k\binom nk = \sum_{k=1}^n n\binom{n-1}{k-1}
= n\sum_{j=0}^{n-1}\binom{n-1}{j} = n\,2^{n-1},
\]
por el \cref{cor:g12:comb:sums}. \emph{Derivando:} al derivar
$(1+x)^n = \sum_k \binom nk x^k$ se obtiene
$n(1+x)^{n-1} = \sum_k k \binom nk x^{k-1}$; se evalúa en $x = 1$.
\end{solution}

\begin{solution}{pb:g12:comb:1}
\textbf{1.} $26^2 \times 10^3 = 676\,000$ matrículas. BANANA: $6$ letras
con la a triplicada y la n duplicada: $\frac{6!}{3!\,2!} = 60$ anagramas.

\textbf{2.} $\binom{32}{5} = 201\,376$ manos;
$\binom42 \binom{28}{3} = 6 \times 3\,276 = 19\,656$ con exactamente dos
ases.

\textbf{3.} Un camino es una palabra con $4$ letras D y $3$ letras A: se
eligen las posiciones de las A: $\binom73 = 35$.

\textbf{4.} $(1+x)^4 = 1 + 4x + 6x^2 + 4x^3 + x^4$. En $x = 1$:
$\sum_k \binom nk = 2^n$; en $x = -1$: $\sum_k (-1)^k \binom nk = 0$: las
sumas de las filas y las sumas alternadas de las filas del \omterm{prop:g11:binom:pascal}{triángulo de
Pascal}.

\textbf{5.} Comisiones de $k$ personas con presidente, elegidas entre $n$:
o bien se elige la comisión y después su presidente
($\binom nk \times k$), o bien el presidente y después los demás miembros
($n \times \binom{n-1}{k-1}$): son iguales. Sumando sobre $k$, el segundo
miembro suma $n \sum_j \binom{n-1}{j} = n\,2^{n-1}$.

\textbf{6.} Una fila de $10$ estrellas y $3$ barras codifica el pedido (las
bolas del sabor 1 antes de la primera barra, etc.); la fila tiene $13$
símbolos y queda determinada por las posiciones de las barras:
$\binom{13}{3} = 286$ pedidos.

\textbf{7.} $12$ estrellas y $2$ barras: $\binom{14}{2} = 91$.

\textbf{8.} Con $x', y', z' \geq 0$ y $x' + y' + z' = 9$:
$\binom{11}{2} = 55$.

\textbf{9.} Un monomio $a^i b^j c^k$ con $i + j + k = 5$:
$\binom72 = 21$.

\textbf{10.} Con la fórmula: $3$ estrellas y $1$ barra:
$\binom41 = 4$; enumerando: $(3,0)$, $(2,1)$, $(1,2)$, $(0,3)$: coinciden.

\textbf{11.} Lo de ``idénticas'' entró cuando se declaró que un pedido no
era más que el \emph{recuento} por sabor: las estrellas no llevan nombre.
Si las bolas se comen en orden, cada una de las $10$ posiciones distintas
elige sabor libremente: $4^{10} = 1\,048\,576$ secuencias; otro modelo y
otro mundo (\cref{met:g12:comb:model}: pregunta siempre
\emph{¿ordenado?, ¿distintos?, ¿con repetición?}).

\textbf{12.} $D_1 = 0$; $D_2 = 1$ (el intercambio); $D_3 = 2$ (los dos
ciclos de longitud $3$); $D_4 = 9$.

\textbf{13.} El invitado 1 recibe el sombrero $k \neq 1$: $n - 1$ opciones.
Si el invitado $k$ recibe el sombrero 1, los $n - 2$ invitados restantes
desarreglan sus propios sombreros: $D_{n-2}$ maneras. Si el invitado $k$
\emph{no} recibe el sombrero 1, renombramos el sombrero 1 como el sombrero
prohibido del invitado $k$: los $n - 1$ invitados restantes desarreglan:
$D_{n-1}$ maneras. Por tanto,
$D_n = (n-1)(D_{n-1} + D_{n-2})$. Comprobación:
$D_4 = 3(2 + 1) = 9$; y $D_5 = 4(9 + 2) = 44$.

\textbf{14.} De los $3! = 6$ repartos, restamos los que fijan al menos un
sombrero: tres fijan un sombrero dado ($2!$ cada uno,
$3 \times 2 = 6$), lo que cuenta de más las parejas ($3$ parejas, $1!$ cada
una), que hay que devolver, y vuelve a restarse la identidad ($1$):
$D_3 = 6 - 6 + 3 - 1 = 2$, es decir,
$3!\left(1 - 1 + \frac12 - \frac16\right) = 2$. En general,
$D_n = n!\sum_{k=0}^{n} \frac{(-1)^k}{k!}$.

\textbf{15.} $\frac{D_5}{120} = \frac{44}{120} \approx 0.3667$, ya cerca de
$\frac1\eu \approx 0.3679$: la suma alternada
$1 - 1 + \frac{1}{2!} - \frac{1}{3!} + \dots$ marcha hacia $\eu^{-1}$. Los
sombreros de una fiesta grande se desarreglan alrededor del $36.8\,\%$ de
las veces: la constante de la lotería y de la secretaria, en su tercera
aparición.

\textbf{16.} $\P(\text{válido}) = \frac{D_{10}}{10!} \approx 0.368$. Cada
repetición del sorteo sale bien con \omterm{def:g10:proba:distribution}{probabilidad} $\approx \frac1\eu$, así
que el número esperado de sorteos es de en torno a $\eu \approx 2.7$: hay
que contar con tres rondas de sombrero.

\textbf{17.} Cada apretón aporta $2$ al recuento total de manos
estrechadas, así que la suma de los números de todos los invitados es par.
Una suma de enteros es par solo si el número de sumandos \omterm{def:g11:func:parity}{impares} es par:
los que estrechan un número \omterm{def:g11:func:parity}{impar} de manos vienen en cantidad par. (Con
tres invitados: ningún perfil de apretones tiene exactamente una o tres
entradas \omterm{def:g11:func:parity}{impares}; compruébalo con los cuatro grafos posibles.)

\textbf{18.} $n = 1$: $1 = 1$. Si
$1^3 + \dots + n^3 = \left(\frac{n(n+1)}{2}\right)^2$, al añadir
$(n+1)^3$:
\[
\frac{n^2(n+1)^2}{4} + (n+1)^3
= \frac{(n+1)^2\left(n^2 + 4n + 4\right)}{4}
= \left(\frac{(n+1)(n+2)}{2}\right)^{\!2} :
\]
paso de inducción hecho. Para $n = 3$: $1 + 8 + 27 = 36 = 6^2$.

\textbf{19.} Un camino hasta $(n, n)$ da $2n$ pasos y cruza la antidiagonal
$x + y = n$ exactamente en un punto de la retícula $(j, n - j)$; la primera
mitad es un camino con $j$ letras D entre $n$ pasos ($\binom nj$
elecciones), y la segunda mitad, leída al revés, otro tanto ($\binom nj$ de
nuevo, por simetría). Sumando sobre el punto de cruce:
$\binom{2n}{n} = \sum_j \binom nj^2$: la identidad de Vandermonde,
dibujada.

\textbf{20.} Multiplicar etapas y sumar casos: las matrículas y las manos
de póquer. Codificar: los caminos como palabras de D y A, y los pedidos
como estrellas y barras. Contar dos veces: la comisión con presidente, los
apretones de manos y los caminos cortados por la mitad. Restar y corregir:
los sombreros desarreglados, con $\frac1\eu$ como residuo. Próximas
paradas: estos recuentos bajo las fracciones de la \omterm{def:g10:proba:distribution}{probabilidad} y el
recuento de caminos con las potencias de las matrices de adyacencia, dos
capítulos más adelante.
\end{solution}
